%%%%%%%%%%%%%%%%%%%%%%%%%%%%%%%%%%%%%%%%%
% Thesis 
% LaTeX Template
% Version 1.0 (01/01/13)
%
% This template has been downloaded from:
% http://www.leonstudio.org
%
% Original authors:
% Leo yan 
%
% License:
% CC BY-NC-SA 3.0 (http://creativecommons.org/licenses/by-nc-sa/3.0/)
%
% Note:
% Make sure to edit document variables in the Thesis.cls file
%
%%%%%%%%%%%%%%%%%%%%%%%%%%%%%%%%%%%%%%%%%

%----------------------------------------------------------------------------------------
%       PACKAGES AND OTHER DOCUMENT CONFIGURATIONS
%----------------------------------------------------------------------------------------
% Paper size, default font size and one-sided paper
\documentclass[11pt, a4paper, oneside]{Thesis}
\usepackage{titlesec}

\usepackage{enumerate}

% Specifies the directory where pictures are stored
\graphicspath{{./Pictures/}}

% Use the natbib reference package - read up on this to edit the reference style; if you want text (e.g. Smith et al., 2012) for the in-text references (instead of numbers), remove 'numbers'
\usepackage[square, numbers, comma, sort&compress]{natbib} 
\hypersetup{urlcolor=blue, colorlinks=true} % Colors hyperlinks in blue - change to black if annoying
\title{\ttitle} % Defines the thesis title - don't touch this


\begin{document}


% Use roman page numbering style (i, ii, iii, iv...) for the pre-content pages
\frontmatter 

% Line spacing of 1.3
\setstretch{1.3}

% Define the page headers using the FancyHdr package and set up for one-sided printing
\fancyhead{} % Clears all page headers and footers
\rhead{\thepage} % Sets the right side header to show the page number
\lhead{} % Clears the left side page header

\pagestyle{fancy} % Finally, use the "fancy" page style to implement the FancyHdr headers

\newcommand{\HRule}{\rule{\linewidth}{0.5mm}} % New command to make the lines in the title page

% PDF meta-data
\hypersetup{pdftitle={\ttitle}}
\hypersetup{pdfsubject=\subjectname}
\hypersetup{pdfauthor=\authornames}
\hypersetup{pdfkeywords=\keywordnames}

%----------------------------------------------------------------------------------------
%       TITLE PAGE
%----------------------------------------------------------------------------------------

\begin{titlepage}
\begin{center}

\includegraphics{logo.jpg}\\[1cm] %
\textsc{\LARGE \univname}\\[1.5cm] % University name
%\textsc{\Large Doctoral Thesis}\\[0.5cm] % Thesis type
\HRule \\[1.0cm] % Horizontal line
{\huge \bfseries \ttitle}\\[0.4cm] % Project

\huge \textit {Project II}\\[0.1cm] % Project
\HRule \\[1.5cm] % Horizontal line

\small \textit{Writen by: \authornames}\\[0.5cm] % Research group name and department name
\textit{Computer Science}\\[2cm] % Research group name and department nam 
\textit{\large Last modified: \today}\\[4cm] % Date

 
\vfill
\end{center}

\end{titlepage}


\clearpage % Start a new page

%----------------------------------------------------------------------------------------
%       QUOTATION PAGE
%----------------------------------------------------------------------------------------

\pagestyle{empty} % No headers or footers for the following pages

\textit{Project Requirement}

\begin{enumerate}
\item 
In this project, implement and/or compare three value numbering techniques, local value numbering
(LVN), super local value numbering (SVN) and dominator-tree-based global value numbering (GVN).
\item GVN has already been implemented in LLVM (see (llvm)/lib/Transforms/Scalar/GVN.cpp).
Note that GVN subsumes SVN, which, in turn, subsumes LVN. You need to update the GVN code so that
it can be downgraded to perform LVN or SVN only. To implement super local value numbering, CS5130
students need to implement an analysis pass that identifies extended basic blocks.
\item You will need to introduce two new command line options for opt which suggest the optimizer to perform
LVN or SVN instead of GVN. I will use command line option -std-compile-opts to trigger standard
1compiler optimizations which include global value numbering. You should add the following two options
to opt:
\\1. -lvn: perform local value numbering only.
\\2. -svn: perform super local value numbering only.
\\3. output EBBs for each function if this option is enabled.
\\
\item Your code must work on the Rekhi lab machines. It is fine to use your own machine for development. But
you must test your code in the lab machines.\\
Download and install test-suite-3.4 from the LLVM website. We will use
test-suite-3.4/SingleSource/Benchmarks/Stanford as the testing benchmark suite. Your
optimizations should pass all the benchmarks in this suite. You also need to collect the total compile time,
the compile time of LVN, SVN and GVN, the total execution time, and the number of instructions deleted
by value numbering.
\item 
Extended Basic Block \\
Dump all the extended basic blocks in each function led by the function name followed by a colon. Each
EBB should be output as a pre-order sequence of the basic blocks in it enclosed by a pair of curly brackets
and each basic block can be represented by its name.
\item
Submission\\
Tar all the source files you have changed including their directories into a tar ball, project2.tar, and submit
it through Canvas. Write a report that summarizes your implementation and analyzes your experimental
results for the Stanford benchmark suite, attach your EBB outputs and name them after the original input file names
with .ebb extension. For example, the EBB output for Bubblesort.c should be stored in file Bubblesort.ebb. 
\end{enumerate}

\clearpage % Start a new page



%----------------------------------------------------------------------------------------
%       LIST OF CONTENTS/FIGURES/TABLES PAGES
%----------------------------------------------------------------------------------------

%\pagestyle{fancy} % The page style headers have been "empty" all this time, now use the "fancy" headers as defined before to bring them back

%\lhead{\emph{Contents}} % Set the left side page header to "Contents"
%\tableofcontents % Write out the Table of Contents



%----------------------------------------------------------------------------------------
%       THESIS CONTENT - CHAPTERS
%----------------------------------------------------------------------------------------

\mainmatter % Begin numeric (1,2,3...) page numbering

\pagestyle{fancy} % Return the page headers back to the "fancy" style

% Include the chapters of the thesis as separate files from the Chapters folder
% Uncomment the lines as you write the chapters

% Chapter 1

%\titleformat{\chapter}{\centering\Huge\bfseries}{Part \,\thechapter\}{1em}{}

%\chapter{Introduction} % Main chapter title

\label{Chapter1} % For referencing the chapter elsewhere, use \ref{Chapter1} 

\lhead{ \emph{Self Spatial Reuse}} % This is for the header on each page - perhaps a shortened title

%----------------------------------------------------------------------------------------

\section{Introduction}

1. compile  the source file

clang -O1 -emit-llvm -c loops.c 

2. run tests   

opt -analyze -ds --ssr=true --cls=64 loops.bc

\section{Implementation}
\subsection{Basic Algorithm}
A memory reference in a loop nest has self-spatial reuse if the
distance (stride) between two consecutive accesses by this reference
in the innermost loop is smaller than the cache line size.

I only think about the self spatial reuse at innermost level ,
 for example I J K order, I only think about
 A[I][J][K], A[I][K][K], A[K][K][K].
If no K at all, I take it as a self-temp and pass over the dectect.
For example,
\begin{lstlisting}
  for (i = 1; i < N; i++)
      for (j = 1; j < N; j++)
          for (k = 1; j < N; j++)
          {
            a[k][j][k] = a[i][j][j] + a[i][k][k]+a[i][i][k];
          }
\end{lstlisting}
the stride of a[i][j][k] is 1, the stride of a[i][k2][k3] is
upperbound of L3 * stride of k2 + stride of k3, k2 means k in level
2 and k3 means k in level 3. Same to a[k1][j][k3]. But when meets
a[i][j][j], I will take it as a self-temp. 

At last we will make sure the stride count will be lower the
cachelinecount, which means two consecutive value store in  a same
cache line.

\subsection{Loop Nest Number} 
Find the outermost loop (no parent) of each loopnest, save it to a
currentLoop, then , find another outermost, check if it is same as the
currentLoop, if same, do nothing, not same, let loopNest++, and make
currentLoop = current outermost loop.

\subsection{Get the variable type} 
Need to know the value type of the array, because different type comes
with different layout size. Int would be 4, and doulbe is 8, which
means a cacheline with size 64 can hold 16 integers but 8 doulbe variable. 
This is one thing we need to notice, if we get a pointer value, we
need to find the element it points to.

\subsection{Get the subscripts} 
   Maily using the Pair struct offered by this DA file. For example 
\begin{lstlisting}
  for (i = 1; i < N; i++)
      for (j = 1; j < N; j++)
      {
        a[i][j] = a[i][2*j] + 10;
      }

\end{lstlisting}
we will get [i] of a[i][j] from Pair[1], and could get the i of a[i]
from Loop 1, also we could get the coeffient 1, which we could use it
as the stride Same idea with [j], we know its coeff is 1 and loop
level is 2, we need its stride is lower or equal to cachelinecount 

\subsection{Get the upperbound}
Upperbound is necessary for a[i][k][k] and a[k][j][k], this kind of
situations. We also need to get all level upperbound in advance.

\subsection{MIV} 
For array A[i][j+k][i], we need do a loop for all j and k in j+k, and
check if it is same as k, if same, we need to use upperbound to
caculate the stride count. 


%----------------------------------------------------------------------------------------
\section{The difficulties}
For A[i][j][j], there is a situation that, cacheline is very big, and
upperbound of lever 2 and 3 is small, then although level 3 has a
self-temp, level 2 still could have a self-reuse, which means
a[i][j][j] and a[i][j+1][j+1] could be in a same cache line.
Considering its big possibily accroding to dimmentions, I did not
handle this case this time. 

%----------------------------------------------------------------------------------------

\newpage
\section{Results}
\begin{lstlisting}
test1

test2

Loop Nest ： 1
  %1 = load i32* %arrayidx, align 8, !tbaa !1
level : 1
stride : 2

  store i32 %1, i32* %arrayidx2, align 4, !tbaa !1
level : 1
stride : 1

test3

Loop Nest ： 1
  %1 = load i32* %arrayidx, align 8, !tbaa !1
level : 1
stride : 4

  store i32 %1, i32* %arrayidx2, align 4, !tbaa !1
level : 1
stride : 2

test4

Loop Nest ： 1
  %2 = load i32* %arrayidx, align 4, !tbaa !1
level : 1
stride : 8

  %4 = load i32* %arrayidx3, align 16, !tbaa !1
level : 1
stride : 12

test5

test6

Loop Nest ： 1
  %3 = load i32* %arrayidx5, align 4, !tbaa !1
level : 2
stride : 1

  store i32 %add10, i32* %arrayidx14, align 4, !tbaa !1
level : 2
stride : 1

test7

Loop Nest ： 1
  %4 = load i32* %arrayidx6, align 8, !tbaa !1
level : 2
stride : 2

  %6 = load i32* %arrayidx11, align 4, !tbaa !1
level : 2
stride : 3

  store i32 %add12, i32* %arrayidx17, align 4, !tbaa !1
level : 2
stride : 1

test8

Loop Nest ： 1
  store i32 %add, i32* %arrayidx13, align 4, !tbaa !1
level : 2
stride : 3

test9

Loop Nest ： 1
  %0 = load double* %arrayidx10, align 8, !tbaa !1
level : 3
stride : 1

  %3 = load double* %arrayidx17, align 8, !tbaa !1
level : 3
stride : 2

  store double %add18, double* %arrayidx25, align 8, !tbaa !1
level : 3
stride : 1

test10

Loop Nest ： 1
  %4 = load double* %arrayidx17, align 8, !tbaa !1
level : 3
stride : 4

  store double %add18, double* %arrayidx24, align 16, !tbaa !1
level : 3
stride : 2

test11

test12

Loop Nest ： 1
  %1 = load double* %arrayidx10, align 8, !tbaa !1
level : 3
stride : 5
\end{lstlisting}

%----------------------------------------------------------------------------------------


\begin{flushright}
%Guide written by ---\\
%Sunil Patel: \href{http://www.sunilpatel.co.uk}{www.sunilpatel.co.uk}
\end{flushright}


%----------------------------------------------------------------------------------------
%       THESIS CONTENT - APPENDICES
%----------------------------------------------------------------------------------------

%\addtocontents{toc}{\vspace{2em}} % Add a gap in the Contents, for aesthetics

%\appendix % Cue to tell LaTeX that the following 'chapters' are Appendices

% Include the appendices of the thesis as separate files from the Appendices folder
% Uncomment the lines as you write the Appendices
%\input{./Appendices/code.pdf}
%\input{./Appendices/AppendixB}
%\input{./Appendices/AppendixC}

\addtocontents{toc}{\vspace{2em}} % Add a gap in the Contents, for aesthetics

\backmatter

\end{document}  